\begin{frame}[t]
\slidehead{предобработка · пайплайн}{Единый слой очистки запроса}
\begin{columns}[T]
\column{0.52\textwidth}
{\InterSemi\fontsize{8.5}{10.5}\selectfont Шесть шагов очистки}\par\vspace{5pt}
{\renewcommand{\arraystretch}{1.22}\setlength{\tabcolsep}{5pt}
\begin{tabular}{|L{2.4cm}|L{4.2cm}|}
\hline
\theadl{Шаг} & \theadl{Что делает}\\
\hline
\tcelll{\hb{Базовая чистка}} & \tcelll{убирает мусорные пробелы, ё $\to$ е}\\
\hline
\tcelll{\hb{Расклейка}} & \tcelll{«128гб» $\to$ «128 гб»}\\
\hline
\tcelll{\hb{Сепараторы}} & \tcelll{«g-pro» $\to$ «g pro»}\\
\hline
\tcelll{\hb{Нормализация}} & \tcelll{единый ключ для словарей}\\
\hline
\tcelll{\hb{Защита брендов}} & \tcelll{«Красный Октябрь» — не цвет}\\
\hline
\tcelll{\hb{Подсказки}} & \tcelll{спаны моделей поверх разметки}\\
\hline
\end{tabular}}
\column{0.46\textwidth}
{\InterSemi\fontsize{8.5}{10.5}\selectfont Зачем один общий слой}\par\vspace{5pt}
{\smalltext
\numdot{1}\ Раньше каждый ноутбук чистил запрос по-своему — результаты не сходились. Теперь очистка одна на всех.\par\vspace{4pt}
\numdot{2}\ Регистр сохраняем там, где он несёт смысл: «красный» — это цвет, а «Красный Октябрь» — бренд конфет.\par\vspace{4pt}
\numdot{3}\ Модуль отдаёт и чистый текст, и подсказки: где в запросе модель, а где защищённый бренд.}
\par\vspace{8pt}
\pinknote{ЭФФЕКТ}{Спан модели появляется у 18\% запросов, склейки чинятся у 11\%, а каждый четвёртый «пустой» хвост после бренда оживает.}
\end{columns}
\slidefoot
\end{frame}

\begin{frame}[t]
\slidehead{предобработка · пример}{Живой пример: запрос до и после}
\graynote{ИСХОДНЫЙ ЗАПРОС}{«Наушники Logitech G-Pro X SE 128гб» — дефисы, склейки, смесь языков.}
\vspace{3pt}\par
{\InterSemi\fontsize{8}{10}\selectfont Шаг 1 · расклейка и сепараторы}\par\vspace{2pt}
{\fontsize{8.5}{12}\selectfont
\wplain{наушники} \ \wplain{logitech} \ \wplain{g} \ \wplain{pro} \ \wplain{x} \ \wplain{se} \ \wplain{128} \ \wplain{гб}}\par\vspace{1pt}
{\smalltext\color{mvgray} «G-Pro» распался в «g pro», «128гб» — в «128 гб». Теперь словари найдут каждую часть.}
\par\vspace{5pt}
{\InterSemi\fontsize{8}{10}\selectfont Шаг 2 · разметка со словарями и подсказками}\par\vspace{2pt}
{\fontsize{8.5}{12}\selectfont
\wtag{mvred}{наушники · КАТЕГОРИЯ} \ \wtag{mvdark}{logitech · БРЕНД} \ \wtag{okgreen}{g pro x se · МОДЕЛЬ} \ \wtag{warnorange}{128 гб · АТРИБУТ}}\par\vspace{1pt}
{\smalltext\color{mvgray} «g pro x se» поймал словарь моделей из майнинга, «128 гб» — регулярка памяти.}
\par\vspace{5pt}
{\InterSemi\fontsize{8}{10}\selectfont Шаг 3 · готовые факты}\par\vspace{2pt}
{\smalltext бренд = Logitech · категория = наушники · модель = g pro x se · память = 128 ГБ.}
\slidefoot
\end{frame}

\begin{frame}[t]
\slidehead{предобработка · измерение}{Что даёт предобработка на семпле}
\begin{columns}[T]
\column{0.5\textwidth}
\vspace{0pt}\centering\includegraphics[width=0.94\textwidth,height=3.7cm,keepaspectratio]{02_preprocess_effect.png}
\column{0.47\textwidth}
\graynote{КАК МЕРИЛИ}{Взяли 12 000 уникальных запросов и прогнали через пайплайн дважды: со словарём моделей и без него.}
\vspace{3pt}\par
{\fontsize{7.2}{9.4}\selectfont
\numdot{1}\ 18\% запросов получают спан модели — примерно каждый пятый.\par\vspace{2pt}
\numdot{2}\ 11\% содержат склейку вида «128гб», которую чинит расклейка.\par\vspace{2pt}
\numdot{3}\ 24\% пустых хвостов после бренда получают тег модели после подсказок.\par\vspace{2pt}
\numdot{4}\ Защита брендов срабатывает у 4{,}5\% — редко, но спасает «Красный Октябрь» от превращения в цвет.}
\end{columns}
\slidefoot
\end{frame}

\begin{frame}[t]
\slidehead{тег model · проблема}{Хвост после бренда теряется}
\begin{columns}[T]
\column{0.44\textwidth}
\vspace{2pt}\centering\includegraphics[width=0.96\textwidth,height=5.9cm,keepaspectratio]{model_tag/01_missing_model_problem.png}
\column{0.53\textwidth}
\pinknote{ПРОБЛЕМА}{После бренда часто идёт название линейки — «dyson v15», «galaxy s24». При слабой разметке эти слова остаются без тега, и CRF не может их выучить.}
\vspace{5pt}\par
{\smalltext
\numdot{1}\ Из 15 тысяч запросов бренд есть у 7046. Почти у половины (46\%) после бренда остался неразмеченный латинский хвост.\par\vspace{3pt}
\numdot{2}\ Типичная длина хвоста — два токена: «zenbook a16», «ipad air».\par\vspace{3pt}
\numdot{3}\ Это не атрибут — нет числа с единицей. Нужен отдельный тег MODEL и словарь линеек.}
\end{columns}
\slidefoot
\end{frame}

\begin{frame}[t]
\slidehead{тег model · примеры}{Как выглядят потерянные хвосты}
\graynote{ДО СЛОВАРЯ}{Серые слова — то, что разметка не смогла объяснить. Именно они и есть модели.}
\vspace{7pt}\par
\begin{columns}[T]
\column{0.48\textwidth}
{\InterSemi\fontsize{8}{10}\selectfont До: хвост без тега}\par\vspace{5pt}
{\fontsize{8.2}{14}\selectfont
\wtag{mvred}{пылесос} \ \wtag{mvdark}{dyson} \ \wplain{v15}\par\vspace{4pt}
\wtag{mvred}{телефон} \ \wtag{mvdark}{samsung} \ \wplain{galaxy} \ \wplain{s24}\par\vspace{4pt}
\wtag{mvdark}{asus} \ \wplain{zenbook} \ \wplain{a16}\par\vspace{4pt}
\wtag{mvdark}{apple} \ \wplain{ipad} \ \wplain{air}}
\column{0.5\textwidth}
{\InterSemi\fontsize{8}{10}\selectfont После: словарь линеек подключён}\par\vspace{5pt}
{\fontsize{8.2}{14}\selectfont
\wtag{mvred}{пылесос} \ \wtag{mvdark}{dyson} \ \wtag{okgreen}{v15 · МОДЕЛЬ}\par\vspace{4pt}
\wtag{mvred}{телефон} \ \wtag{mvdark}{samsung} \ \wtag{okgreen}{galaxy s24 · МОДЕЛЬ}\par\vspace{4pt}
\wtag{mvdark}{asus} \ \wtag{okgreen}{zenbook a16 · МОДЕЛЬ}\par\vspace{4pt}
\wtag{mvdark}{apple} \ \wtag{okgreen}{ipad air · МОДЕЛЬ}}
\end{columns}
\vspace{7pt}\par
\pinknote{ЧТО ИЗМЕНИЛОСЬ}{Хвост после бренда перестал быть «мусором» — теперь это полноценный факт «модель», по которому можно искать и ранжировать.}
\slidefoot
\end{frame}

\begin{frame}[t]
\slidehead{тег model · майнинг}{Майнинг словаря: пример «g pro»}
\graynote{ИДЕЯ}{Названия кликнутых карточек уже содержат модели. Находим в названии бренд — и всё, что идёт после него, считаем кандидатом в модели.}
\vspace{4pt}\par
\begin{columns}[T]
\column{0.49\textwidth}
{\InterSemi\fontsize{8}{10}\selectfont Название карточки $\to$ кандидаты}\par\vspace{3pt}
{\fontsize{6.8}{10}\selectfont
\wplain{Наушники} \wtag{mvdark}{Logitech} \wtag{okgreen}{G} \wtag{okgreen}{PRO} \wtag{okgreen}{X} \wtag{okgreen}{SE} \wplain{Black}}\par\vspace{3pt}
{\smalltext После бренда снимаем 1–4 коротких токена:}\par\vspace{2pt}
{\renewcommand{\arraystretch}{1.05}\setlength{\tabcolsep}{4pt}
\fontsize{7}{9}\selectfont
\begin{tabular}{|L{2.0cm}|C{1.9cm}|}
\hline
\theadl{Кандидат} & \thead{Раз}\\
\hline
\tcelll{g} & \tcellc{213}\\
\hline
\tcelll{\hb{g pro}} & \tcellc{\hnum{87}}\\
\hline
\tcelll{g pro x} & \tcellc{31}\\
\hline
\tcelll{g pro x se} & \tcellc{9}\\
\hline
\tcelll{g pro x se black} & \tcellc{2 — отсекли}\\
\hline
\end{tabular}}
\column{0.48\textwidth}
{\InterSemi\fontsize{8}{10}\selectfont Почему порог — 6 раз}\par\vspace{3pt}
{\smalltext
\numdot{1}\ Фраза попадает в словарь, если встретилась \textbf{минимум 6 раз}. «g pro» видели 87 раз — берём. «g pro x se black» — два раза, случайный хвост, отсекаем.\par\vspace{3pt}
\numdot{2}\ Порог 2 — словарь раздувается до 7215 фраз и тащит мусор. Порог 20 — остаётся 335 фраз, теряются редкие линейки.\par\vspace{3pt}
\numdot{3}\ Шесть — золотая середина: 1512 фраз, 22\% хвостов покрыто, словарь чист на 96{,}6\%.}
\end{columns}
\slidefoot
\end{frame}

\begin{frame}[t]
\slidehead{тег model · порог}{Выбор порога: размер словаря против покрытия}
\begin{columns}[T]
\column{0.5\textwidth}
\vspace{0pt}\centering\includegraphics[width=0.94\textwidth,height=3.8cm,keepaspectratio]{model_tag/03_threshold_tradeoff.png}
\column{0.47\textwidth}
\graynote{КАК ЧИТАТЬ ГРАФИК}{Чёрная кривая — размер словаря, красная — доля хвостов, которые он покрывает. Пунктир — наш выбор.}
\vspace{3pt}\par
{\fontsize{7.2}{9.4}\selectfont
\numdot{1}\ Из 47 тысяч названий карточек намайнили 8540 сырых кандидатов: galaxy, iphone, watch, redmi и другие.\par\vspace{2.5pt}
\numdot{2}\ Кривые расходятся плавно, «локтя» нет — выбрали порог 6 как компромисс между размером и покрытием.\par\vspace{2.5pt}
\numdot{3}\ Дополнительные фильтры: голые числа не берём; фразы из одного токена — только коды вида «v15», «a305», «ps5».}
\end{columns}
\slidefoot
\end{frame}

\begin{frame}[t]
\slidehead{тег model · результат}{До и после подключения словаря}
\begin{columns}[T]
\column{0.5\textwidth}
\vspace{2pt}\centering\includegraphics[width=0.98\textwidth,height=4.4cm,keepaspectratio]{model_tag/04_before_after_model_tag.png}
\column{0.47\textwidth}
\conclusion{Словарь линеек включает класс MODEL: было 0 запросов с тегом — стало 1573. Потерянных хвостов стало меньше на четверть.}
\vspace{6pt}\par
{\smalltext
\numdot{1}\ Латинские хвосты без тега: 3262 $\to$ 2502, минус 23\%.\par\vspace{4pt}
\numdot{2}\ Проверили качество словаря вручную: 96{,}6\% фраз — настоящие линейки, остальное — длинные артикулы и аксессуары, их вычистим.\par\vspace{4pt}
\numdot{3}\ Оставшиеся хвосты словарём не закрыть — это кандидаты на ручную разметку, для неё мы сделали приложение (покажем дальше).}
\end{columns}
\slidefoot
\end{frame}

\begin{frame}[t]
\slidehead{разметка · схемы}{Какие бывают схемы токен-разметки}
\graynote{КОНТЕКСТ}{Задача NER — найти спаны сущностей. Но модели классифицируют токены, поэтому спаны кодируют в токен-схему. Разные схемы — разные способы обозначить границы.}
\vspace{3pt}\par
{\renewcommand{\arraystretch}{1.12}\setlength{\tabcolsep}{5pt}
\begin{tabular}{|L{2.0cm}|C{2.3cm}|L{7.3cm}|}
\hline
\theadl{Схема} & \thead{Теги} & \theadl{Суть}\\
\hline
\tcelll{\hb{IO}} & \tcellc{I, O} & \tcelll{самое простое, но две соседние сущности сливаются}\\
\hline
\tcelll{\hb{BIO (IOB)}} & \tcellc{B, I, O} & \tcelll{B отмечает начало — границы видны; стандарт CoNLL}\\
\hline
\tcelll{\hb{IOE}} & \tcellc{I, O, E} & \tcelll{то же самое, только отмечается конец, а не начало}\\
\hline
\tcelll{\hb{BILOU}} & \tcellc{B, I, L, O, U} & \tcelll{добавляются теги конца (L) и одиночного слова (U)}\\
\hline
\tcelll{\hb{BIES}} & \tcellc{B, I, E, S, O} & \tcelll{максимум классов — сильнее дисбаланс, учить труднее}\\
\hline
\end{tabular}}
\vspace{3pt}\par
\tnote{Общая беда всех токен-схем: вложенные сущности не выразить — «адрес» нельзя разбить на город, улицу и дом одновременно.}
\slidefoot
\end{frame}

\begin{frame}[t]
\slidehead{разметка · выбор}{Храним спаны, учим на BIO}
\conclusion{Спаны — универсальный формат хранения и оценки: любой модельный ответ приводится к ним. BIO — формат обучения. Держим оба.}
\vspace{6pt}\par
\begin{columns}[T]
\column{0.5\textwidth}
{\InterSemi\fontsize{8.5}{10.5}\selectfont Почему именно BIO}\par\vspace{4pt}
{\smalltext
\numdot{1}\ Классы — полные теги (B-BRAND, I-ATTR и так далее). Буквы B, I, O лишь кодируют границы.\par\vspace{3pt}
\numdot{2}\ У 27\% запросов сущности длиннее одного токена. Схема IO склеила бы соседние сущности в одну, BIO их разделяет.\par\vspace{3pt}
\numdot{3}\ BILOU на коротких запросах даёт больше классов без выигрыша — дисбаланс растёт, качество нет.\par\vspace{3pt}
\numdot{4}\ Наши теги: BRAND, MODEL, CATEGORY, ATTR (с приставками B и I) плюс O.}
\column{0.47\textwidth}
\begin{tikzpicture}
\node[fill=cardbg, rounded corners=3pt, inner sep=8pt, text width=\dimexpr\textwidth-16pt\relax]{
{\InterSemi\fontsize{8}{10}\selectfont Пример BIO-разметки}\par\vspace{5pt}
{\fontsize{8}{12}\selectfont\color{mvdark}
ноутбук \ \biotag{mvred}{B-CATEGORY}\par\vspace{2pt}
asus \ \biotag{mvdark}{B-BRAND}\par\vspace{2pt}
zenbook \ \biotag{okgreen}{B-MODEL}\par\vspace{2pt}
16 \ \biotag{warnorange}{B-ATTR}\par\vspace{2pt}
гб \ \biotag{warnorange}{I-ATTR}}
};
\end{tikzpicture}
\end{columns}
\slidefoot
\end{frame}

\begin{frame}[t]
\slidehead{mvp · архитектура}{Гибридный каскад: MVP на тот день}
\pinknote{КАК УСТРОЕНО}{К концу дня 03 каскад выглядел так. Каждый слой закрывает слабость предыдущего: правила мгновенные, CRF понимает морфологию, дальше три классификатора — бренд, категория и тип атрибута. На выходе факты за десятки миллисекунд.}
\vspace{5pt}\par
\begin{center}
\begin{tikzpicture}[node distance=0pt]
  \node (a) {\pipestepn{mvred}{white}{ЗАПРОС\\\fontsize{5}{6}\selectfont «ноутбук asus 16 гб»}};
  \node[right=1pt of a] (c1) {\chevron};
  \node[right=1pt of c1] (b) {\pipestepn{cardbg}{mvdark}{ПРАВИЛА\\\fontsize{5}{6}\selectfont словари, regex}};
  \node[right=1pt of b] (c2) {\chevron};
  \node[right=1pt of c2] (d) {\pipestepn{cardbg}{mvdark}{CRF\\\fontsize{5}{6}\selectfont разметка токенов}};
  \node[right=1pt of d] (c3) {\chevron};
  \node[right=1pt of c3] (e) {\pipestepn{cardbg}{mvdark}{БРЕНД\\\fontsize{5}{6}\selectfont классификатор}};
  \node[right=1pt of e] (c4) {\chevron};
  \node[right=1pt of c4] (f) {\pipestepn{cardbg}{mvdark}{КАТЕГОРИЯ\\\fontsize{5}{6}\selectfont классификатор}};
  \node[right=1pt of f] (c5) {\chevron};
  \node[right=1pt of c5] (g) {\pipestepn{mvdark}{white}{АТРИБУТЫ\\\fontsize{5}{6}\selectfont классификатор}};
\end{tikzpicture}
\end{center}
\vspace{5pt}\par
\begin{center}
\begin{tikzpicture}
  \node[fill=pinkband, rounded corners=3pt, inner xsep=10pt, inner ysep=6pt,
        text width=\dimexpr\textwidth-40pt\relax, align=left] {%
    {\color{mvred}\InterBold\fontsize{8}{10}\selectfont ПЕТЛЯ КАЧЕСТВА}\hspace{7pt}%
    {\color{mvdark}\fontsize{8}{11}\selectfont неуверенные ответы и ошибки уходят в ручную разметку, на ней дообучаем CRF и пополняем словари}};
\end{tikzpicture}
\end{center}
\slidefoot
\end{frame}

\begin{frame}[t]
\slidehead{продакшен · путь запроса}{Путь запроса от строки до выдачи}
\graynote{ЗАЧЕМ}{Извлечь факты — только середина пути: дальше нормализация, каталог и ранжирование.}
\vspace{5pt}\par
\begin{center}
\begin{tikzpicture}[node distance=0pt]
  \node (p1) {\pipestepq{cardbg}{mvdark}{Запрос}};
  \node[right=1pt of p1] (a1) {\chevron};
  \node[right=1pt of a1] (p2) {\pipestepq{cardbg}{mvdark}{Пред-\\обработка}};
  \node[right=1pt of p2] (a2) {\chevron};
  \node[right=1pt of a2] (p3) {\pipestepq{mvred}{white}{Извлечение\\фактов}};
  \node[right=1pt of p3] (a3) {\chevron};
  \node[right=1pt of a3] (p4) {\pipestepq{cardbg}{mvdark}{Нормали-\\зация}};
  \node[right=1pt of p4] (a4) {\chevron};
  \node[right=1pt of a4] (p5) {\pipestepq{cardbg}{mvdark}{Каталог}};
  \node[right=1pt of p5] (a5) {\chevron};
  \node[right=1pt of a5] (p6) {\pipestepq{cardbg}{mvdark}{Ранжиро-\\вание}};
\end{tikzpicture}
\end{center}
\vspace{5pt}\par
{\renewcommand{\arraystretch}{1.1}\setlength{\tabcolsep}{7pt}
\begin{center}
\begin{tabular}{|L{3.2cm}|L{8.4cm}|}
\hline
\theadl{Шаг} & \theadl{Пример}\\
\hline
\tcelll{\hb{Нормализация}} & \tcelll{«16гб», «16 gb», «16 ГБ» — всё превращается в одно значение: 16 ГБ}\\
\hline
\tcelll{\hb{Каталог}} & \tcelll{категория «ноутбук» находит узел «Ноутбуки и планшеты»}\\
\hline
\tcelll{\hb{Ранжирование}} & \tcelll{карточки с брендом asus и памятью 16 ГБ идут наверх}\\
\hline
\end{tabular}
\end{center}}
\vspace{-2pt}
\slidefoot
\end{frame}

\begin{frame}[t]
\slidehead{mvp · классификатор бренда}{Silver для бренда: 67 классов, не только бренды}
\pinknote{КЛЮЧЕВАЯ ИДЕЯ}{Классификатор — запасной слой, когда словарь не нашёл бренд: на «холодильник» нельзя тыкать в частый клик-бренд. Поэтому в классы добавили NO\_BRAND и UNKNOWN.}
\vspace{2pt}\par
\begin{columns}[T]
\column{0.46\textwidth}
\centering\includegraphics[width=0.95\textwidth,height=3.9cm,keepaspectratio]{brand_clf/mv_02_silver_composition.png}
\column{0.51\textwidth}
{\fontsize{7}{9}\selectfont
\numdot{1}\ Из 63 тысяч агрегированных запросов в silver попали 28 664: у остальных нет надёжного сигнала — их честно выбросили.\par\vspace{1.5pt}
\numdot{2}\ Уверенность метки — из кликов: клик выше в выдаче весит больше, майор-бренд должен набрать не меньше 55\% веса.\par\vspace{1.5pt}
\numdot{3}\ NO\_BRAND — запрос-категория без бренда. UNKNOWN — бренд есть, но вне топ-65 («стинол», «darina»): модель учится отказываться, а не угадывать похожий.\par\vspace{1.5pt}
\numdot{4}\ Разбиение: 22 846 на обучение, 5 710 на валидацию.}
\end{columns}
\slidefoot
\end{frame}

\begin{frame}[t]
\slidehead{mvp · классификатор бренда}{Четыре модели: что и как обучали}
\graynote{ПОСТАНОВКА}{Мультикласс запрос $\to$ бренд на silver. Все четыре — линейные модели на TF-IDF: отвечают за миллисекунды, критично для SLA. Разница — в признаках и балансировке.}
\vspace{3pt}\par
{\renewcommand{\arraystretch}{0.95}\setlength{\tabcolsep}{4pt}
\fontsize{6.6}{8.2}\selectfont
\begin{tabular}{|L{2.6cm}|L{4.2cm}|L{4.2cm}|}
\hline
\theadl{Модель} & \theadl{Признаки} & \theadl{Особенность}\\
\hline
\tcelll{\hb{LogReg слова+симв.}} & \tcelll{слова 1–2 + символьные n-граммы 2–4} & \tcelll{выбрана в прод: меньше всего ложных брендов}\\
\hline
\tcelll{SGD символьная} & \tcelll{символьные n-граммы} & \tcelll{лучший macro-F1, но чаще врёт на категориях}\\
\hline
\tcelll{LogReg символьная} & \tcelll{символьные n-граммы} & \tcelll{базовый вариант без словных признаков}\\
\hline
\tcelll{LogReg с балансом} & \tcelll{символьные n-граммы} & \tcelll{веса классов — хуже на редких брендах}\\
\hline
\end{tabular}}
\par\vspace{3pt}
{\smalltext
\numdot{1}\ Символьные кусочки по 2–4 буквы делают «айфон», «afon» и «iphone» похожими — опечатки не страшны.\par\vspace{2pt}
\numdot{2}\ Веса по уверенности клика: сомнительные метки влияют меньше. Все модели сохранены, сравнение — на одной валидации.}
\slidefoot
\end{frame}

\begin{frame}[t]
\slidehead{mvp · классификатор бренда}{Результаты: 0,95 macro-F1 и честные отказы}
\begin{columns}[T]
\column{0.46\textwidth}
\vspace{0pt}\centering\includegraphics[width=0.94\textwidth,height=4.6cm,keepaspectratio]{brand_clf/mv_01_models_compare.png}
\column{0.51\textwidth}
\graynote{КАК ЧИТАТЬ}{Верх — macro-F1 до и после порогов отказа. Низ — доля запросов-категорий, где модель ошибочно навязала бренд (главная метрика, цель — ноль).}
\vspace{3pt}\par
{\fontsize{7.1}{9.2}\selectfont
\numdot{1}\ По macro-F1 SGD формально впереди (0{,}951 против 0{,}950), но разница — третий знак.\par\vspace{2pt}
\numdot{2}\ Решает вторая метрика: у LogReg слова+символы ложный бренд на категориях 6{,}9\% против 9{,}4\% у SGD — её и берём в прод.\par\vspace{2pt}
\numdot{3}\ F1 спец-классов: NO\_BRAND 0{,}88, UNKNOWN 0{,}86 — модель научилась говорить «бренда нет» и «не знаю».\par\vspace{2pt}
\numdot{4}\ Примеры: «айфон 15 про» даёт Apple, «холодильник» — NO\_BRAND.}
\end{columns}
\slidefoot
\end{frame}

\begin{frame}[t]
\slidehead{mvp · классификатор бренда}{Пороги отказа: лучше пусто, чем ложный бренд}
\begin{columns}[T]
\column{0.48\textwidth}
\vspace{0pt}\centering\includegraphics[width=0.94\textwidth,height=3.4cm,keepaspectratio]{brand_clf/mv_04_thresholds.png}
\par\vspace{2pt}
{\fontsize{6.8}{8.6}\selectfont\color{mvgray} Хвост по классам: PlayStation 0{,}67 (4 примера), UNKNOWN 0{,}86, NO\_BRAND 0{,}88 — редкие классы ждут золотой разметки.}
\column{0.49\textwidth}
{\fontsize{7}{9}\selectfont
{\InterSemi Когда классификатор вообще зовут:}\par\vspace{2pt}
\numdot{1}\ Словарь или алиас уже нашли бренд — не зовут.\par\vspace{2pt}
\numdot{2}\ Запрос — чистая категория — не зовут, бренд пустой.\par\vspace{2pt}
\numdot{3}\ Есть намёк на бренд (алиас, похожее на модель слово) — зовут.\par\vspace{3pt}
{\InterSemi Как принимается ответ:}\par\vspace{2pt}
\numdot{4}\ Уверенность top-1 не ниже 0{,}42 и отрыв от top-2 не меньше 0{,}08 — принимаем бренд.\par\vspace{2pt}
\numdot{5}\ Ответ NO\_BRAND или UNKNOWN выше своего порога — бренд остаётся пустым, и это правильно.\par\vspace{2pt}
\numdot{6}\ Пороги не прошли — отказ, бренд пустой. С порогами точность даже растёт: отказы съедают в основном ошибки.}
\end{columns}
\slidefoot
\end{frame}

\begin{frame}[t]
\slidehead{приложение · разметка}{Приложение разметки на C++/Qt: мастер из трёх этапов}
\graynote{КАК УСТРОЕНО}{Нативное приложение на C++ и Qt/QML. Запрос разбит на блоки по центру, фон с историей заблюрен. Этап 1: клавиши B/I/O, метка появляется над блоком, Backspace снимает и шагает назад, Enter — подтверждение этапа.}
\vspace{5pt}\par
\begin{columns}[T]
\column{0.49\textwidth}
\centering
\appshot{width=\dimexpr\textwidth-8pt\relax,height=4.6cm,keepaspectratio}{apps/cpp_label_stage1.png}\par\vspace{2pt}
{\smalltext\color{mvgray} Этап 1 — разметка B/I/O: активный блок в красной рамке, метки над блоками}
\column{0.49\textwidth}
\centering
\appshot{width=\dimexpr\textwidth-8pt\relax,height=4.6cm,keepaspectratio}{apps/cpp_label_stage2.png}\par\vspace{2pt}
{\smalltext\color{mvgray} Этап 2 — тип сущности: у каждого типа описание, когда его ставить}
\end{columns}
\slidefoot
\end{frame}

\begin{frame}[t]
\slidehead{приложение · умный поиск}{MVP умного поиска на C++/Qt}
\pinknote{КИЛЛЕР-ФИЧА}{Ранжирование каталога опирается только на извлечённые факты — бренд, категорию, модель и атрибуты, — а не на сырой текст запроса. Движок на C++ отвечает за 1–2 миллисекунды.}
\vspace{5pt}\par
\begin{columns}[T]
\column{0.49\textwidth}
\centering
\appshot{width=\dimexpr\textwidth-8pt\relax,height=4.6cm,keepaspectratio}{apps/cpp_mvp_search.png}\par\vspace{2pt}
{\smalltext\color{mvgray} Факты чипами + товары, отранжированные по фактам, с полосой релевантности}
\column{0.49\textwidth}
\centering
\appshot{width=\dimexpr\textwidth-8pt\relax,height=4.6cm,keepaspectratio}{apps/cpp_mvp_json.png}\par\vspace{2pt}
{\smalltext\color{mvgray} По клику раскрывается JSON с сущностями; таймер извлечения — в шапке}
\end{columns}
\slidefoot
\end{frame}
